%% kapa: exam=1 round=34 year=2023 period=2 form=A subject=appraisal_law subject_ko=감정평가관계법규 seq=12 number=12
%% source: Q-Net 감정평가사 기출문제 (cst00309) artlSeq=5211789 file=raw/1008/5211789/2023년도 제34회 감정평가사 1차시험 2교시.pdf
%% conversion: kordoc 4.24.51 -> tools/gen1.py (auto); review=proofread-sonnet; flags=none
\begin{question}{12}
국토의 계획 및 이용에 관한 법령상 도시·군계획시설사업에 관한 설명으로 옳은 것은?
\begin{choices}
  \item 대도시 시장이 작성한 도시·군계획시설사업에 관한 실시계획은 국토교통부장관의 인가를 받아야 한다.
  \item 도시·군계획시설사업이 둘 이상의 시 또는 군의 관할 구역에 걸쳐 시행되게 되는 경우에는 국토교통부장관이 시행자를 정한다.
  \item 도시·군계획시설사업의 대상시설을 둘 이상으로 분할하여 도시·군계획시설사업을 시행할 수 없다.
  \item 「한국토지주택공사법」에 따른 한국토지주택공사가 도시·군계획시설사업의 시행자로 지정받기 위해서 제출해야 하는 신청서에 자금조달계획은 포함되지 않는다.
  \item 「한국전력공사법」에 따른 한국전력공사는 도시·군계획시설사업의 시행자가 될 수 있다.
\end{choices}
\end{question}
\answer{5}
